% ============================================================================
% ARBEITSBLATT: Wiederholung Mi casa
% ============================================================================
% Sequenz 3, Block d: Wiederholung vor Stundenleistung
% 20 Punkte gesamt
% ============================================================================

\documentclass[11pt, a4paper]{article}

% ============================================================================
% SW-MODUS TOGGLE
% ============================================================================
\newif\ifswmodus
\swmodusfalse

% ============================================================================
% PAKETE
% ============================================================================

\usepackage[utf8]{inputenc}
\usepackage[T1]{fontenc}
\usepackage[ngerman]{babel}
\usepackage{helvet}
\renewcommand{\familydefault}{\sfdefault}

\usepackage[margin=2cm]{geometry}
\usepackage{enumitem}
\usepackage{tabularx}
\usepackage{booktabs}
\usepackage{multirow}
\usepackage{pgffor}

\usepackage{xcolor}
\usepackage{framed}
\usepackage{tcolorbox}
\tcbuselibrary{skins}

\usepackage{fancyhdr}
\usepackage{lastpage}
\usepackage{needspace}

% ============================================================================
% FARBEN
% ============================================================================

\definecolor{spanisch-color}{HTML}{c41e3a}
\definecolor{grau-color}{HTML}{888888}
\definecolor{hellgrau-color}{HTML}{f5f5f5}
\definecolor{orange-color}{HTML}{b35c00}

\definecolor{sw-dark}{HTML}{000000}
\definecolor{sw-gray}{HTML}{666666}
\definecolor{sw-white}{HTML}{ffffff}

\ifswmodus
    \colorlet{spanisch}{sw-dark}
    \colorlet{grau}{sw-gray}
    \colorlet{hellgrau}{sw-white}
    \colorlet{orange}{sw-dark}
\else
    \colorlet{spanisch}{spanisch-color}
    \colorlet{grau}{grau-color}
    \colorlet{hellgrau}{hellgrau-color}
    \colorlet{orange}{orange-color}
\fi

\newcommand{\fachfarbe}{spanisch}

% ============================================================================
% VARIABLEN
% ============================================================================

\newcommand{\thema}{Wiederholung: Mi casa}
\newcommand{\sequenz}{03}
\newcommand{\block}{d}
\newcommand{\fach}{Spanisch}
\newcommand{\klasse}{7}
\newcommand{\maxpunkte}{20}
\newcommand{\datum}{\today}

% ============================================================================
% KOPF- UND FUSSZEILE
% ============================================================================

\pagestyle{fancy}
\fancyhf{}
\renewcommand{\headrulewidth}{0.4pt}
\renewcommand{\footrulewidth}{0.4pt}

\lhead{\fach{} -- Klasse \klasse}
\chead{AB-\sequenz\block}
\rhead{\datum}

\lfoot{}
\cfoot{}
\rfoot{Seite \thepage\ von \pageref{LastPage}}

% ============================================================================
% BEFEHLE
% ============================================================================

\newcommand{\punktebox}[1]{\hfill\fbox{\small \_/#1}}
\newcommand{\luecke}[1]{\rule{#1}{0.4pt}}

\newtcolorbox{merkkasten}[1][]{
    colback=hellgrau,
    colframe=\fachfarbe,
    colbacktitle=white,
    coltitle=\fachfarbe,
    boxrule=1pt,
    arc=0pt,
    left=8pt, right=8pt, top=8pt, bottom=8pt,
    fonttitle=\bfseries,
    attach boxed title to top left={yshift=-2mm, xshift=5mm},
    boxed title style={boxrule=0pt, colback=white},
    title=#1
}

\newtcolorbox{vokabelkasten}{
    colback=white,
    colframe=\fachfarbe,
    boxrule=0.5pt,
    arc=0pt,
    left=5pt, right=5pt, top=5pt, bottom=5pt
}

\newtcolorbox{hinweis}{
    colback=white,
    colframe=orange,
    colbacktitle=white,
    coltitle=orange,
    boxrule=1pt,
    arc=0pt,
    left=5pt, right=5pt, top=5pt, bottom=5pt,
    fonttitle=\bfseries,
    attach boxed title to top left={yshift=-2mm, xshift=5mm},
    boxed title style={boxrule=0pt, colback=white},
    title=Hinweis
}

\newenvironment{aufgabe}[1]{%
    \needspace{5\baselineskip}%
    \nopagebreak[4]%
    \vspace{10pt}
    \noindent\textbf{\textcolor{\fachfarbe}{Aufgabe #1}}
    \nopagebreak[4]%
    \vspace{5pt}
    \nopagebreak[4]%
    \begin{list}{}{\leftmargin=0pt}
    \item[]
}{%
    \end{list}
}

\newlist{teilaufgaben}{enumerate}{2}
\setlist[teilaufgaben,1]{label=\alph*), leftmargin=2em}
\setlist[teilaufgaben,2]{label=\roman*), leftmargin=2em}

\newcommand{\antwortzeilen}[1]{%
    \vspace{2pt}
    \foreach \n in {1,...,#1}{%
        \noindent\rule{\textwidth}{0.4pt}\vspace{8pt}
    }
}

\newcommand{\es}[1]{\textcolor{\fachfarbe}{\textbf{#1}}}

% ============================================================================
% DOKUMENT
% ============================================================================

\begin{document}

% ============================================================================
% KOPFBEREICH
% ============================================================================

\begin{center}
    {\Large\bfseries\textcolor{\fachfarbe}{Arbeitsblatt: \thema}}
\end{center}

\vspace{5pt}

\noindent
\begin{tabularx}{\textwidth}{|X|X|X|}
\hline
\textbf{Name:} \luecke{4cm} &
\textbf{Klasse:} \klasse &
\textbf{Datum:} \luecke{2cm} \\
\hline
\end{tabularx}

\vspace{10pt}

\begin{hinweis}
Dieses Arbeitsblatt hilft dir, fur die Stundenleistung zu uben. Wiederhole alle Themen der Sequenz ,,Mi casa''!
\end{hinweis}

\vspace{10pt}

% ============================================================================
% MERKKASTEN: Raume
% ============================================================================

\begin{merkkasten}[Merkkasten: Las habitaciones (Die Raume)]
\textbf{Ubertrage die Vokabeln aus dem Lernpfad:}

\vspace{0.5em}

\begin{tabular}{ll}
\es{la cocina} & = \luecke{4cm} \\[0.8em]
\es{el salon} & = \luecke{4cm} \\[0.8em]
\es{el dormitorio} & = \luecke{4cm} \\[0.8em]
\es{el bano} & = \luecke{4cm} \\[0.8em]
\es{el jardin} & = \luecke{4cm} \\
\end{tabular}
\end{merkkasten}

\vspace{10pt}

% ============================================================================
% AUFGABE 1: Raume zuordnen
% ============================================================================

\begin{aufgabe}{1}
Verbinde die spanischen Raume mit der deutschen Ubersetzung. \punktebox{4}
\end{aufgabe}

\begin{center}
\begin{tabular}{rcl}
\es{la cocina} & \luecke{3cm} & das Schlafzimmer \\[0.8em]
\es{el salon} & \luecke{3cm} & der Garten \\[0.8em]
\es{el dormitorio} & \luecke{3cm} & die Kuche \\[0.8em]
\es{el jardin} & \luecke{3cm} & das Wohnzimmer \\
\end{tabular}
\end{center}

\vspace{10pt}

% ============================================================================
% AUFGABE 2: hay + Mobel
% ============================================================================

\begin{aufgabe}{2}
Erganze die Satze mit \es{hay} und dem passenden Mobel. \punktebox{4}
\end{aufgabe}

\begin{vokabelkasten}
\textit{Mobel:} una cama | un sofa | una mesa | un armario
\end{vokabelkasten}

\vspace{0.5em}

\noindent
a) En el dormitorio \luecke{4cm} \luecke{3cm}. \\[0.8em]
b) En el salon \luecke{4cm} \luecke{3cm}. \\[0.8em]
c) En la cocina \luecke{4cm} \luecke{3cm}. \\[0.8em]
d) En el dormitorio \luecke{4cm} \luecke{3cm}. \\

\vspace{10pt}

% ============================================================================
% AUFGABE 3: Farben + Adjektivanpassung
% ============================================================================

\begin{aufgabe}{3}
Passe die Farbadjektive an das Nomen an. \punktebox{6}
\end{aufgabe}

\begin{vokabelkasten}
\textit{Farben:} rojo/roja | azul | verde | blanco/blanca | negro/negra | amarillo/amarilla
\end{vokabelkasten}

\vspace{0.5em}

\noindent
a) La mesa es \luecke{3cm} (rot). \\[0.8em]
b) El sofa es \luecke{3cm} (blau). \\[0.8em]
c) La silla es \luecke{3cm} (gelb). \\[0.8em]
d) El armario es \luecke{3cm} (weiß). \\[0.8em]
e) Las camas son \luecke{3cm} (grun). \\[0.8em]
f) Los sofas son \luecke{3cm} (schwarz). \\

\vspace{10pt}

% ============================================================================
% AUFGABE 4: Zimmer beschreiben
% ============================================================================

\begin{aufgabe}{4}
Beschreibe dein Zimmer auf Spanisch. Verwende mindestens 4 Satze mit \es{hay} und Farbadjektiven. \punktebox{6}
\end{aufgabe}

\begin{hinweis}
Beispiel: En mi dormitorio hay una cama blanca. Hay un armario grande...
\end{hinweis}

\vspace{0.5em}

\noindent\textbf{Mi dormitorio:}

\antwortzeilen{6}

% ============================================================================
% PUNKTEUBERSICHT
% ============================================================================

\vspace{20pt}
\hrule
\vspace{10pt}

\begin{flushright}
\textbf{Punkte:} \luecke{1cm} / \maxpunkte
\end{flushright}

\end{document}
